% results.tex -- Experiments and Discussion

\section{Experiments}
\label{sec:experiments}

\subsection{Evaluation Protocol}

We propose the following protocol for evaluating an AI agent on
\textsc{MachiaveliBench}:

\begin{enumerate}[nosep]
    \item \textbf{Full tournament}: run three episodes for every
          (game, strategy) pair, yielding a total of 306 episodes.
    \item \textbf{Metric computation}: aggregate per-episode results
          into the six metrics defined in Section~\ref{sec:metrics}.
    \item \textbf{Reporting}: present per-game breakdowns alongside
          composite scores to enable fine-grained analysis of
          behavioural patterns.
\end{enumerate}

\subsection{Baseline Agents}

To illustrate the benchmark's discriminative power, we outline three
baseline agent types:

\begin{description}[nosep]
    \item[Random Agent.] Selects uniformly from available actions each
        round. This agent serves as a lower bound and is expected to
        exhibit near-uniform cooperation rates with no adaptability.
    \item[Tit-for-Tat Agent.] Cooperates on the first round and then
        mirrors the opponent's previous action. This classic strategy
        achieves high cooperation with reciprocal partners but is
        exploitable in the first round by adversarial opponents.
    \item[LLM Agent.] A language-model agent that receives the
        observation as a structured prompt and selects an action via
        natural-language reasoning. The LLM agent's behaviour depends
        on its pretraining distribution and instruction-tuning,
        making it the primary subject of interest for alignment
        research.
\end{description}

\subsection{Expected Behavioural Patterns}

Based on game-theoretic analysis and prior
work~\cite{axelrod1984evolution,pan2023machiavelli}, we anticipate the
following patterns:

\begin{itemize}[nosep]
    \item \textbf{Cooperation--defection trade-off}: agents optimizing
          purely for reward will defect in the Prisoner's Dilemma but
          may cooperate in the Stag Hunt where mutual cooperation is
          a Nash equilibrium~\cite{nash1950equilibrium}.
    \item \textbf{Fairness norms}: instruction-tuned LLMs are expected
          to propose near-equal splits in the Ultimatum Game, reflecting
          human-like fairness preferences instilled during
          training~\cite{christiano2017deep}.
    \item \textbf{Trust calibration}: well-calibrated agents should
          invest moderately in the Trust Game and adjust returns based
          on the trustee's observed generosity.
    \item \textbf{Free-riding incentives}: in the Public Goods Game,
          purely self-interested agents will contribute minimally,
          while agents with social preferences will contribute more.
\end{itemize}

\section{Discussion}
\label{sec:discussion}

\subsection{Implications for Alignment}

The metrics defined in \textsc{MachiaveliBench} are directly relevant
to AI safety concerns identified by Amodei et
al.~\cite{amodei2016concrete}:

\begin{itemize}[nosep]
    \item \textbf{Reward hacking}: an agent that achieves high scores
          by consistently defecting against cooperative opponents
          exposes a misalignment between reward maximization and
          socially desirable behaviour.
    \item \textbf{Safe exploration}: the exploitation resistance
          metric quantifies how well an agent avoids being exploited
          by adversarial partners, a prerequisite for safe deployment
          in open environments.
    \item \textbf{Scalable oversight}: fairness and Pareto efficiency
          provide proxy signals for whether an agent's strategy
          produces outcomes that a human supervisor would endorse.
\end{itemize}

\subsection{Emergence of Cooperation}

Axelrod's tournaments~\cite{axelrod1984evolution} demonstrated that
reciprocal strategies outperform exploitative ones in iterated
settings. \textsc{MachiaveliBench} tests whether modern AI agents
rediscover this result. An agent that develops conditional cooperation---high
cooperation with reciprocal partners and measured defection against
exploiters---demonstrates sophisticated strategic reasoning without
requiring explicit game-theoretic programming.

\subsection{Trust Calibration}

The Trust Game and Ultimatum Game probe whether agents can calibrate
their behaviour to the norms of their interaction partner. An
alignment-oriented agent should be willing to trust generous partners
more than stingy ones and should reject unfair offers in the Ultimatum
Game to enforce fairness norms. The adaptability metric captures this
calibration by measuring variance in cooperation rate across different
opponent types.

\subsection{Limitations}

Several limitations should guide interpretation of benchmark results:

\begin{itemize}[nosep]
    \item The opponent pool, while diverse, does not capture the full
          spectrum of adversarial or cooperative behaviours encountered
          in real-world multi-agent systems.
    \item Sequential games are single-shot, limiting the study of
          reputation effects and long-term trust dynamics.
    \item The benchmark evaluates dyadic interactions; extending to
          $n$-player settings would increase ecological validity.
    \item Payoff structures are fixed; adaptive payoff schedules could
          test robustness to distributional shift.
\end{itemize}
